% Part C -- Knowledge quiz, FORM B (20 items).
%
% Item Bnn occupies the same blueprint cell as item Ann in part_c_form_a.tex:
% same topic, same Bloom level, same target difficulty, same targeted
% misconception -- different scenario, no shared stem wording.
% See docs/blueprint.md for the full mapping.
%
% Key positions are balanced: 5 items each with the key at (a), (b), (c), (d).
%
% The quiz was cut from 28 to 20 items on 2026-08-05. The eight retired A/B
% pairs are archived verbatim in common/deprecated_items.tex, which also records
% the old-to-new numbering.

\partheading{C}{Knowledge questions}

\instruction{20 questions, each with exactly one best answer. Tick one box per question. If you do not
know, please tick ``I do not know'' rather than guessing --- an honest blank is more useful to us than a
lucky guess, and none of this affects your grade. Plan for about 15 minutes.}

\begin{questions}

% ===== The model and its context =====================================

% --- Item 1: T1 / remember / easy -------------------------------------
\qitem Which statement about the LLM inside an agentic system is correct?
\begin{choices}
  \choice It accumulates knowledge about the current task across the calls of one run.
  \choice It keeps the tool results it has already seen, so they need to be sent only once.
  \CorrectChoice It has no memory of previous calls, so the harness has to resend everything it needs.
  \choice It remembers the system prompt after the first call and does not require it again.
\end{choices}
\dontknow
\itemmeta{T1 --- LLM role \& statelessness}{remember}{easy}{B01}
\rationale{All three distractors are forms of the same error --- assuming server-side state --- and each is
a bug students actually write when assembling requests.}

% --- Item 2: T2 / understand / medium ---------------------------------
\qitem A team observes that their agent follows instructions reliably in short runs but starts ignoring
them in long ones, although the context window is never exceeded. What is the most plausible explanation?
\begin{choices}
  \choice The system prompt is silently truncated once the conversation passes a certain number of messages.
  \choice Instructions expire after a fixed number of turns and have to be repeated.
  \choice The model switches to a smaller variant automatically to keep response times constant.
  \CorrectChoice The model's effective attention to any single instruction falls as the window fills up with lower-signal material.
\end{choices}
\dontknow
\itemmeta{T2 --- Context engineering}{understand}{medium}{B02}
\rationale{Same cell as A02: degradation begins well before the hard limit. The distractors invent
mechanisms that would produce an abrupt rather than a gradual failure.}

% --- Item 3: T2 / apply-analyse / hard --------------------------------
\qitem Your harness summarises the oldest third of the conversation whenever the context exceeds a
threshold. What are you accepting in exchange?
\begin{choices}
  \choice The agent can no longer call any tool that was used in the summarised part of the run.
  \CorrectChoice Detail from the summarised turns is no longer available in full, so some fidelity is traded for sustained coherence and lower cost.
  \choice The summary has to be produced by a different model from the one running the loop.
  \choice The loop's termination conditions stop applying, because the iteration count is reset.
\end{choices}
\dontknow
\itemmeta{T2 --- Context engineering}{apply-analyse}{hard}{B03}
\rationale{Same cell as A03, approached from the cost side rather than the choice side: compaction is a
deliberate loss of fidelity, and knowing what you gave up is the point.}

% ===== The loop, tools, MCP and skills ===============================

% --- Item 4: T3 / understand / medium ---------------------------------
\qitem Your agent occasionally gets stuck calling the same failing tool over and over until you kill the
process. Which property of the loop is missing?
\begin{choices}
  \choice A retry policy that repeats each failing tool call until it succeeds.
  \choice A larger context window, so that the repeated observations still fit.
  \choice A sandbox, so that repeated tool calls cannot damage the host system.
  \CorrectChoice A bound on execution --- for example a maximum number of iterations --- that forces the machine into a terminal state.
\end{choices}
\dontknow
\itemmeta{T3 --- Loop and control flow}{understand}{medium}{B04}
\rationale{Same cell as A04. (a) is the opposite of what is needed and (c) attributes termination to the
sandbox.}

% --- Item 5: T3 / apply-analyse / hard --------------------------------
\qitem Some harnesses build planning instructions into the system prompt rather than waiting for a user
to ask for a plan. What follows directly from that design?
\begin{choices}
  \CorrectChoice Every run spends its first reasoning effort on producing a plan, whether or not the task needs one.
  \choice The user can no longer influence how the agent approaches the task.
  \choice The agent's tool calls are validated against the plan by the tool registry.
  \choice The plan replaces the conversation history as the agent's context.
\end{choices}
\dontknow
\itemmeta{T3 --- Loop and control flow}{apply-analyse}{hard}{B05}
\rationale{Same cell as A05: up-front planning is a trade, not a free improvement --- it costs the first
turn of every run, including trivial ones.}

% --- Item 6: T4 / remember / easy -------------------------------------
\qitem Which of these would be implemented as a tool rather than as some other component of the harness?
\begin{choices}
  \choice A bundle of instructions that teaches the agent how to fill in a particular PDF form.
  \choice A callback that runs before every action to check whether it is permitted.
  \CorrectChoice A function that lists the files in a directory and returns the listing to the model.
  \choice A policy that decides whether an action requires the user's confirmation.
\end{choices}
\dontknow
\itemmeta{T4 --- Tools and tool design}{remember}{easy}{B06}
\rationale{Same cell as A06, asked as a classification instead of a definition: (a) is a skill, (b) a hook,
(d) permission management.}

% --- Item 7: T4 / apply-analyse / hard --------------------------------
\qitem Your harness catches every exception raised inside a tool and returns the literal string
\texttt{error} to the model. Which consequence follows most directly?
\begin{choices}
  \choice The harness can no longer log the original exception, because it has been replaced.
  \CorrectChoice The model has nothing to work with in order to choose a corrective action, so it cannot recover from failures it could otherwise handle.
  \choice The tool's schema becomes invalid, because its declared return type is not met.
  \choice The agent treats every failure as a final answer and ends the run.
\end{choices}
\dontknow
\itemmeta{T4 --- Tools and tool design}{apply-analyse}{hard}{B07}
\rationale{Same cell as A07 with the ambiguity introduced deliberately rather than accidentally: errors
must reach the model in a form it can act on. Logging (a) is unaffected.}

% --- Item 8: T5 / remember / easy -------------------------------------
\qitem In the architecture of a harness, what role does an MCP server play?
\begin{choices}
  \choice It hosts the language model that the agent sends its requests to.
  \choice It stores the agent's conversation history so that it survives a restart.
  \CorrectChoice It provides tools that the agent can discover and call over a standardised interface.
  \choice It enforces which of the agent's actions require the user's approval.
\end{choices}
\dontknow
\itemmeta{T5 --- MCP}{remember}{easy}{B08}
\rationale{Same cell as A08: (a) is the ``MCP serves the model'' confusion, (b) long-term memory,
(d) permission management.}

% --- Item 9: T5 / apply-analyse / medium ------------------------------
\qitem A team adds a new tool to their MCP server. Their agent picks it up on the next run with no code
change and no configuration change. Which property of their integration does that demonstrate?
\begin{choices}
  \CorrectChoice Tools are discovered at runtime, so the agent's action space follows the server's.
  \choice The tool schemas are compiled into the harness at build time.
  \choice The model was fine-tuned on the server's tool catalogue.
  \choice The tool registry caches every tool it has ever seen and replays it.
\end{choices}
\dontknow
\itemmeta{T5 --- MCP}{apply-analyse}{medium}{B09}
\rationale{Same cell as A09, reasoning from an observation back to the mechanism. (b) would produce the
opposite behaviour.}

% --- Item 10: T6 / understand / medium --------------------------------
\qitem A harness holds a library of 60 skills, but a given run typically loads two or three of them.
What mechanism makes that possible?
\begin{choices}
  \CorrectChoice The skills are disclosed just in time: only an identifier and a short description are visible until one is actually needed.
  \choice The skills are compressed and transparently decompressed by the context manager.
  \choice The agent requests skills from the model provider's servers rather than from the harness.
  \choice The unused skills are held in long-term memory and summarised into the context.
\end{choices}
\dontknow
\itemmeta{T6 --- Skills}{understand}{medium}{B10}
\rationale{Same cell as A10: the mechanism is selective disclosure, not compression --- (b) is the common
wrong intuition that the problem is size rather than signal.}

% ===== Memory, hooks, permissions, sandboxing and sub-agents =========

% --- Item 11: T7 / remember / easy ------------------------------------
\qitem Which requirement can be met \emph{only} by a long-term memory component, and not by context
management alone?
\begin{choices}
  \choice Older turns of the current conversation are summarised in order to save tokens.
  \CorrectChoice Information stored during one session is still available after the application has been restarted.
  \choice The tool schemas stay visible to the model throughout the run.
  \choice The agent's reasoning from the previous iteration informs the next one.
\end{choices}
\dontknow
\itemmeta{T7 --- Long-term memory}{remember}{easy}{B11}
\rationale{Same cell as A11: persistence beyond the process is the defining property. (a), (c) and (d) are
all achievable within a single run's context.}

% --- Item 12: T8 / remember / easy ------------------------------------
\qitem At which points would a hook typically be triggered?
\begin{choices}
  \choice Only once, when the agent process starts up.
  \choice Whenever the user types a new message, and at no other time.
  \choice Whenever the context window is about to overflow, and at no other time.
  \CorrectChoice Before a tool call and after an LLM response.
\end{choices}
\dontknow
\itemmeta{T8 --- Hooks}{remember}{easy}{B12}
\rationale{Same cell as A12: hooks sit at the transitions of the loop, which is what makes them the natural
place for enforcement and instrumentation.}

% --- Item 13: T8 / apply-analyse / medium -----------------------------
\qitem Your team has to record a metric for every tool call, every model call and every permission
decision, without scattering counter increments throughout the codebase. Which mechanism gives you that?
\begin{choices}
  \choice A tool that the agent calls whenever it wants something to be counted.
  \choice A wrapper around the metrics endpoint that infers the events from HTTP traffic.
  \CorrectChoice Lifecycle hooks at those points, with the instrumentation attached to the hooks.
  \choice A sub-agent whose task is to watch the main agent and report what it did.
\end{choices}
\dontknow
\itemmeta{T8 --- Hooks}{apply-analyse}{medium}{B13}
\rationale{Same cell as A13, using the observability side of the same seam. (a) makes measurement depend on
the model choosing to measure.}

% --- Item 14: T9 / understand / easy ----------------------------------
\qitem In a model with a read-only tier, a sandbox-edit tier and a full-access tier, which action is
characteristic of the sandbox-edit tier?
\begin{choices}
  \CorrectChoice Writing a file inside the container's workspace.
  \choice Sending an email on the user's behalf.
  \choice Searching the repository for a string.
  \choice Deleting a branch on the remote repository.
\end{choices}
\dontknow
\itemmeta{T9 --- Permission management}{understand}{easy}{B14}
\rationale{Same cell as A14, classifying into the middle tier instead of the top one: (c) is read-only,
while (b) and (d) reach outside the sandbox.}

% --- Item 15: T9 / understand / medium --------------------------------
\qitem Your permission system is configured so that every single tool call requires the user's
confirmation. Which criticism of that configuration is valid?
\begin{choices}
  \choice It makes the permission decisions impossible to record in the observability data.
  \CorrectChoice It gives up the benefit of the tiers: routine, low-risk observation interrupts the user as often as genuinely consequential action.
  \choice It means the sandbox is no longer needed, since the user approves everything.
  \choice It prevents the agent from using tools discovered from an MCP server.
\end{choices}
\dontknow
\itemmeta{T9 --- Permission management}{understand}{medium}{B15}
\rationale{Same cell as A15 approached from over-restriction: tiers exist to spend human attention where
it changes the outcome. Confirmation fatigue makes approvals reflexive and defeats the gate.}

% --- Item 16: T10 / understand / medium -------------------------------
\qitem A team argues that because their agent runs inside a container, they do not need permission
management. Why is that argument wrong?
\begin{choices}
  \CorrectChoice A sandbox bounds what the agent \emph{can} reach; permission management decides which of the reachable actions it \emph{may} take, including ones whose effects leave the sandbox.
  \choice Containers isolate only the network and not the file system, so file access still needs a policy.
  \choice Permission management is what starts and stops the container, so it cannot be removed.
  \choice A container's isolation expires after a fixed time, after which a policy has to take over.
\end{choices}
\dontknow
\itemmeta{T10 --- Sandboxing}{understand}{medium}{B16}
\rationale{Same cell as A16: \emph{can} versus \emph{may}. A containerised agent with a network connection
and an API token can still act irreversibly on the outside world.}

% --- Item 17: T11 / understand / medium -------------------------------
\qitem Why does giving a sub-agent its own context window help the system scale?
\begin{choices}
  \choice The two agents can call the LLM simultaneously, halving the run's latency.
  \choice The sub-agent's context does not count towards the provider's token limits.
  \choice The main agent can then run on a smaller model, since the reasoning happens in the sub-agent.
  \CorrectChoice The coordinating agent's context stays focused, because only the summary of the sub-agent's work is added to it.
\end{choices}
\dontknow
\itemmeta{T11 --- Sub-agents and A2A}{understand}{medium}{B17}
\rationale{Same cell as A17: the benefit is context economy, not parallelism (a) or a billing loophole (b).}

% ===== Observability and quality assurance ===========================

% --- Item 18: T12 / understand / easy ---------------------------------
\qitem A dashboard for an agent harness shows only the number of completed runs per hour. What is the
most important thing it is missing?
\begin{choices}
  \choice The version number of the harness that produced them.
  \choice The number of lines of Python in the tool registry.
  \CorrectChoice Whether those runs succeeded or failed.
  \choice The physical location of the server that ran them.
\end{choices}
\dontknow
\itemmeta{T12 --- Observability and telemetry}{understand}{easy}{B18}
\rationale{Same cell as A18: throughput without an outcome dimension cannot distinguish a healthy system
from one that is failing fast and often.}

% --- Item 19: T12 / apply-analyse / medium ----------------------------
\qitem Why is it not enough to expose only aggregate metrics for an agentic system?
\begin{choices}
  \choice Aggregate metrics cannot be scraped by a Prometheus-compatible server.
  \choice Aggregate metrics are only valid if the agent is deterministic, which it is not.
  \choice Aggregate metrics are computed by the model itself and are therefore unreliable.
  \CorrectChoice They show that something is wrong but not why; reconstructing what happened in one run requires per-run traces or structured logs.
\end{choices}
\dontknow
\itemmeta{T12 --- Observability and telemetry}{apply-analyse}{medium}{B19}
\rationale{Same cell as A19, stated as a principle rather than an incident. This is why the lab asks for
both a dashboard and inspectable per-run records.}

% --- Item 20: T13 / understand / medium -------------------------------
\qitem Why should the assertions of an end-to-end test for an agent target the resulting state rather
than the agent's wording?
\begin{choices}
  \choice Because the agent's messages are not available to the test harness.
  \CorrectChoice Because the same correct outcome can be reached with different phrasings and different sequences of calls.
  \choice Because state assertions are the only kind a test framework is able to express.
  \choice Because the wording of the agent's messages is fixed by the model provider.
\end{choices}
\dontknow
\itemmeta{T13 --- Quality assurance}{understand}{medium}{B20}
\rationale{Same cell as A20 as a principle: asserting on wording produces flaky tests, because the model's
phrasing varies while the outcome does not.}

\end{questions}
